% TOP_PROBLEM: {"id": 3343, "title": "Vertex Cover Integrality Gap", "category_id": 18, "category": {"id": 18, "name": "logic", "display_name": "Logic", "description": "Problems in mathematical logic, model theory, proof theory, and finite model theory.", "slug": "logic", "order_index": 18, "created_at": "2026-07-31T15:26:25.671Z"}, "status": "open", "problem_number": "OPG-38188", "set_id": 12}
% FIRST_POSTED: 2026-10-02
% ATTEMPT_STATUS: unresolved
% MODEL: GPT 6 Astra Ultra
% UnsolvedMath ID 3343; source catalogue code OPG-38188
% !TeX program = lualatex
% Research attempt dated 2026-10-02; canonical statement preserved.
\documentclass[11pt,a4paper]{article}
\usepackage[margin=25mm]{geometry}
\usepackage{fontspec}
\setmainfont{lmroman10-regular.otf}[BoldFont=lmroman10-bold.otf,
 ItalicFont=lmroman10-italic.otf,BoldItalicFont=lmroman10-bolditalic.otf]
\usepackage{amsmath,amssymb,mathtools}
\usepackage{unicode-math}
\setmathfont{latinmodern-math.otf}
\AtBeginDocument{\let\setminus\smallsetminus}
\usepackage{xurl,hyperref}
\hypersetup{colorlinks=true,urlcolor=blue}
\setcounter{secnumdepth}{0}
\setlength{\parindent}{0pt}
\setlength{\parskip}{5pt}
\setlength{\emergencystretch}{3em}
\begin{document}
\section*{Vertex Cover Integrality Gap}\label{3343}
{\small UnsolvedMath ID 3343 · OPG-38188 · Logic · OpenGarden}
\subsection{Definitions and mathematical statement}
All graphs are finite, simple and undirected. A vertex cover of $G$ is a
set of vertices meeting every edge; write $\operatorname{vc}(G)$ for its
minimum possible cardinality. In the graph vocabulary with adjacency and
equality, $C^k$ is first-order logic using at most $k$ distinct variable
names, which can be reused, and counting quantifiers
$\exists^{\ge m}x\,\varphi(x)$ for integer $m\ge0$. This quantifier says
that at least $m$ vertices satisfy $\varphi$. Write $G\equiv^C_k H$ when
$G$ and $H$ satisfy exactly the same $C^k$ sentences, with no restriction
on sentence length or quantifier depth.

The proposed assertion is that for every $0<\varepsilon<1$ there are
$\delta>0$ and $n_0$ such that for every integer $n\ge n_0$ there are
$n$-vertex graphs $G,H$, with $\operatorname{vc}(H)>0$, satisfying
\[
 G\equiv^C_{\lfloor\delta n\rfloor}H,
 \qquad \operatorname{vc}(G)\ge(2-\varepsilon)\operatorname{vc}(H).
\]
Increase $n_0$ if necessary so the variable bound is positive. The
positive-cover condition makes explicit the intended nontrivial ratio:
two edgeless graphs would otherwise satisfy the displayed inequality
vacuously. The constants may depend on $\varepsilon$, but not on $n$.
The target is the exact counting-logic assertion, including its linear
number of variables.
\subsection{Short English statement}
Can two large graphs look identical to counting logic with linearly many variables while their minimum vertex covers differ by a factor arbitrarily close to two?
\subsection{Research attempt}
\paragraph{Scope and method.}
This research attempt by GPT-6 Astra Ultra is dated 2 October 2026.
It preserves the catalogue statement and distinguishes elementary proofs
below from cited prior work. No novelty claim or formal proof-assistant
verification is made. The final paragraph states the precise remaining scope.
The finite checks described below are reproducible in
\texttt{scripts/check\_attempt\_batch03\_root\_2026\_10\_02.py}.
\paragraph{The number of variables is decisive.}
We construct pairs with vertex-cover ratio arbitrarily
close to two that are indistinguishable in the full
two-variable counting logic, at every quantifier depth.
Then we exhibit a three-variable sentence distinguishing
them. This proves a restricted version and pinpoints
why it cannot meet the catalogue's requirement of
$\lfloor\delta n\rfloor$ variables for fixed $\delta>0$.
Atserias's problem discussion [R1] distinguishes the
two-variable case from stronger variable bounds; the
argument below does not turn a low-variable example
into the asserted linear-variable theorem.

\paragraph{A regular pair with a near-two cover gap.}
Fix an integer $d\ge2$ and put $N_0=2d(d+1)$.
Let $G_0$ consist of $2d$ disjoint copies of $K_{d+1}$,
and let $H_0$ consist of $d+1$ disjoint copies of
$K_{d,d}$. Both graphs have $N_0$ vertices and are
$d$-regular. A cover of $K_{d+1}$ needs all but one
vertex and this suffices. A cover of $K_{d,d}$ has
size at least $d$ because of its perfect matching,
and either part attains that bound. Consequently
\[
 \operatorname{vc}(G_0)=2d^2,\quad
 \operatorname{vc}(H_0)=d(d+1),\quad
 \frac{\operatorname{vc}(G_0)}{\operatorname{vc}(H_0)}
       =\frac{2d}{d+1}=2-\frac2{d+1}.
\]
For any $n\ge N_0$, write $n=qN_0+r$ with
$0\le r<N_0$. Take $q$ copies of each base graph
and add $r$ isolated vertices to each. The ratio is
unchanged and the denominator is positive. Choosing
$d$ sufficiently large gives the desired numerical
ratio for any fixed positive $\varepsilon$ and every
sufficiently large order, in the two-variable setting.

\paragraph{Why all two-variable counting sentences agree.}
Partition each graph into its isolated and regular
vertices. For a pair of named vertices, retain only
their two classes and whether they are equal or adjacent.
Induction on formulas shows that the truth of every
two-variable counting formula is determined by this
information, with the same answer in both graphs.
Atomic formulas have that property. Boolean operations
preserve it. When one variable is counted while the
other is fixed at a regular vertex, the available
choices consist of one equal vertex, $d$ neighbours,
$qN_0-d-1$ other regular vertices, and $r$ isolates.
For a fixed isolated vertex there are no neighbours,
one equal isolate, $r-1$ other isolates and $qN_0$
regular vertices. These counts agree in the two graphs.
Thus every counting threshold preserves the induction,
as do ordinary quantifiers. Sentences have no free
vertices left and obtain identical truth values.
The proof applies at arbitrary depth, not only to a
fixed number of refinement rounds.

\paragraph{The exact obstruction to using more variables.}
With three variables one can state that distinct
vertices $x,y,z$ are pairwise adjacent. Our graph $G$
satisfies this triangle sentence because $d+1\ge3$;
$H$ does not, since its nontrivial components are
bipartite. For every fixed positive $\delta$, the
allowed number $\lfloor\delta n\rfloor$ eventually
reaches three. The exhibited pair then fails the
logical equivalence required in the problem.

There is also a useful relaxation calculation. Summing
the fractional vertex-cover edge inequalities on a
$d$-regular graph with $N$ vertices gives
$d\sum_v x_v\ge dN/2$. The all-half assignment
attains $N/2$. Isolates contribute zero. Hence both
graphs have identical fractional optimum, despite
their very different integral covers. Equality of
this elementary relaxation does not certify equivalence
in a counting logic with more variables.

\paragraph{Verification and final status.}
The root script constructs pairs for $2\le d\le6$,
several copy counts and isolated remainders, and checks
the degree classes and triangle distinction exactly.
The component arguments prove the cover values and
the structural induction proves full two-variable
equivalence. Achieving comparable gaps while hiding
them from linearly many variables remains unresolved
by this attempt.

\subsection{Sources}
\begin{itemize}
\item[S1] ULAM.AI and the UnsolvedMath Contributors, supplied problems.json, record 3343 (OPG-38188), OpenGarden collection. Catalogue identity and source transcription; CC BY 4.0.
\url{https://huggingface.co/datasets/ulamai/UnsolvedMath}.
\item[S2] Open Problem Garden, Vertex Cover Integrality Gap. Original problem and discussion; the mathematical formulation below is expanded from the supplied source record.
\url{http://www.openproblemgarden.org/op/vertex_cover_integrality_gap}.

\item[R1] A. Atserias, Vertex-Cover Integrality Gaps in Counting Logic,
Section 4.1 of Dagstuhl Report 12(1), Seminar 22051 (2022).
\url{https://drops.dagstuhl.de/storage/04dagstuhl-reports/volume12/issue01/22051/html/DagRep.12.1.101/DagRep.12.1.101.html}.

\end{itemize}
\end{document}
